\section{Introduction}

Granular materials exhibit mechanical behavior that emerges from interactions between individual particles and the evolving network of contacts through which forces are transmitted. 
Numerical simulation of granular compaction is therefore a challenging computational mechanics problem. 
During compaction, the contact topology can evolve substantially, frictional contact can govern the transmission of stresses, and large deformation may lead to highly localized mechanical response and, eventually, particle-scale failure or fracture. 
These characteristics make the numerical representation of granular assemblies fundamentally different from that of a conventional continuum with a fixed material domain.

Material point methods (MPM) and other particle-based approaches are consequently attractive for many large-deformation granular mechanics problems, since they can accommodate evolving configurations without relying on a conventional mesh that must remain well behaved throughout the deformation. 
Finite element methods (FEM), however, provide a complementary computational framework in which the continuum formulation, constitutive response, contact treatment, boundary conditions, and numerical solution procedure can be specified and investigated systematically. 
Rather than assuming a priori that one numerical method is preferable for all granular mechanics problems, this project uses FEM as a controlled computational framework for investigating how far a programmable finite element formulation can be developed for realistic three-dimensional granular assemblies.

The \textit{Granular RVEs} project therefore develops a three-dimensional computational mechanics framework for realistic granular microstructures using FEniCSx. 
The project begins with well-defined benchmark problems at the individual-grain level and progressively introduces interacting grains, frictional contact, periodic boundary conditions, and realistic three-dimensional grain geometries. 
This progression is intended to establish a verified computational foundation in which the physical and numerical complexity of the simulations can be increased in a controlled manner.

A central motivation for this development is the eventual comparison of finite element and alternative numerical formulations for granular compaction. 
Realistic grain geometries and representative volume element (RVE) configurations provide a natural basis for such comparisons because the same microstructure, material parameters, loading conditions, and boundary conditions can in principle be supplied to different numerical frameworks. 
Once the FEniCSx formulation is sufficiently developed and verified, these matched simulations can support one-to-one comparisons with other approaches, including material point methods, with respect to mechanical response, numerical behavior, computational cost, and practical limitations.

The longer-term objective is to extend the computational framework from direct grain-scale simulation toward homogenization, model parameter identification, and uncertainty quantification. 
These capabilities would provide a basis for connecting resolved microstructural simulations with reduced-order and data-driven models. 
In particular, the resulting parameterized finite element simulations may ultimately serve as the high-fidelity model evaluations required for the development and assessment of surrogate models and neural-operator approaches for granular mechanics.

At its current stage, the project focuses on establishing this computational foundation rather than claiming a complete solution to granular compaction. 
The emphasis is on reproducible scientific software, verification of the underlying numerical formulations, and a progressive increase in physical complexity. 
This provides a framework in which the capabilities and limitations of the FEniCSx approach can be evaluated rigorously as the project develops.

